%% parts/p12/sat-and-certificates.tex --- Part 12.1
\chapter{SAT Search and Certificates}
\label{sec:p12-sat-and-certificates}

\noindent\emph{Reader note.}
This chapter is a separately authored editorial correction and expansion after
the exact inherited Automated Reasoning chapter.  Its residual-state DPLL
presentation and annex proof retain their own labels and evidence standing;
they neither replace nor alter the source's clause-set presentation, literal
pseudocode, theorem, or inline proof.

The canonical clause and CNF definitions are in
Section~\ref{sec:prop-syntax-semantics-04}.  The exact inherited
equisatisfiability distinction and worked Tseitin constraints are in
Section~\ref{sec:automated-reasoning-02}.  The boundary used here is therefore
explicit: a Tseitin output is an expanded-signature equisatisfiable clause set,
never silently promoted to same-signature formula equivalence.

\section{The displayed DPLL logical core}

\begin{hbdefinition}{def:p12-dpll-state}{Residual state and unit propagation}
Let \(F\) be a finite set of finite propositional clauses.  In this setting,
\(\operatorname{Var}(F)\defeq
\bigcup_{C\in F}\operatorname{Var}(C)\); keep this original variable set
fixed throughout the search.  A \emph{residual state} is a
pair \((G,u)\), where \(u:\operatorname{Var}(F)\rightharpoonup\{0,1\}\) is a
consistent partial truth assignment.  For each original clause not already
satisfied by \(u\), its \emph{residual} is obtained by deleting exactly the
literals that \(u\) makes false; \(G\) is the set of all and only those
residuals.  A \emph{unit-propagation step} selects
a singleton clause \(\{L\}\), extends \(u\) so that \(L\) is true, deletes
clauses containing \(L\), and deletes the complementary literal from every
remaining clause.  The residual-state name is necessary to distinguish the
changing search object from the immutable input \(F\); it is metalinguistic,
not a new clause syntax, and the distinction is used in the correctness
statement and annex proof.  Pure-literal elimination is a different
satisfiability-preserving optimization: if an unassigned variable occurs in
the finite residual clause set with only one polarity, assign it the value that
satisfies those occurrences and delete every clause containing them.  It is an
optional SAT simplification, not unit propagation, because its assignment is
licensed by one-sided occurrence rather than forced by a singleton clause.
The distinction is used in Theorem~\ref{thm:p12-dpll-correctness}, whose exact
displayed core excludes pure-literal elimination.
\index{residual state}\index{unit propagation}\index{pure-literal elimination}
\end{hbdefinition}

\begin{hbmethod}{def:p12-dpll}{DPLL logical core}
\index{DPLL}
The identifier \(\mathsf{DPLL}\) is needed for this exact finite-clause search
method; it names only the unit-propagation and binary-branching operations
displayed here.  Starting from
\((F,\varnothing)\), repeatedly unit-propagate.  If no residual clauses remain,
extend \(u\) arbitrarily to every still-unassigned member of
\(\operatorname{Var}(F)\) and return that total valuation.  If the empty clause
appears, report failure of the current branch.  Otherwise choose a previously
unassigned original variable \(p\), search the state obtained by setting
\(p\) true, and, if that branch fails, search the state obtained by setting
\(p\) false.  Thus \(\mathsf{DPLL}\) stands as an algorithmic method over the
Part-2 clause semantics, not as an inference calculus; its exact distinction
from the later conflict-analysis apparatus is used in
Table~\ref{tab:p12-dpll-cdcl}.
\end{hbmethod}

\begin{hbtheorem}{thm:p12-dpll-correctness}{proved=annex:p12.sat.1}{Correctness of the displayed DPLL logical core}
For every finite clause set \(F\), Method~\ref{def:p12-dpll} terminates.  It
returns a total valuation of the fixed set \(\operatorname{Var}(F)\) satisfying
\(F\) if and only if \(F\) is satisfiable, and it reports failure at the root if
and only if \(F\) is unsatisfiable.  This status covers unit propagation and
binary branching exactly as displayed; it does not cover pure-literal
elimination or any added search, certificate, or implementation mechanism.
\end{hbtheorem}

\begin{hbexample}{def:p12-dpll-trace}{A unit-propagation refutation trace}
For
\[
 F=\bigl\{\{p,q\},\{\neg p,r\},\{\neg q,r\},\{\neg r\}\bigr\},
\]
the first unit clause fixes \(r\) false.  The second and third clauses become
the unit clauses \(\{\neg p\}\) and \(\{\neg q\}\).  Propagating them fixes
\(p\) and \(q\) false, turning the first clause into the empty clause.  The
exhausted root search therefore reports unsatisfiable.  This trace checks this
one finite instance; it is not a completeness proof for DPLL.
\end{hbexample}

\section{Certificates and practical status}

\begin{hbdefinition}{def:p12-sat-certificate-boundary}{SAT and UNSAT evidence}
A \emph{SAT certificate} for \(F\) is a total valuation whose direct clause
check accepts.  An \emph{UNSAT certificate} is a finite refutation log in a
declared proof-logging format accepted by its independently specified checker.  These
names are needed because a solver's response token, a model, a checked
refutation, and trust in solver code establish different facts.  They are
metalinguistic evidence objects rather than new semantic truth values, and the
distinction controls the solver-audit method in
  Chapter~\ref{ch:formal-practice}.
\index{SAT certificate}\index{UNSAT certificate}\index{proof logging}
\end{hbdefinition}

Under the fixed effective encoding used in Parts~2 and~4,
\(\mathsf{SAT}\) is NP-complete under polynomial-time many-one reductions
(Theorems~\ref{thm:p02-sat-np-complete} and~\ref{thm:p04-cook-levin};
\cite{Cook1971}).  This does not say that every instance is hard, that DPLL is
exponential on every input, or that \(\mathsf P\ne\mathsf{NP}\).  A
polynomial-time algorithm for all SAT instances would prove
\(\mathsf P=\mathsf{NP}\); success on structured benchmarks proves neither
direction.

Horn-SAT and 2-SAT remain homed at Theorems~\ref{thm:p02-horn-sat} and
\ref{thm:p02-two-sat}; XOR-SAT is Theorem~\ref{thm:p02-xor-sat}.  The master
algorithm-and-complexity comparison is Table~\ref{tab:p02-tractable-sat}, not
a duplicate P12 table.  A \emph{random \(k\)-SAT
distribution} fixes a variable count, a probability law for generating
\(k\)-literal clauses, and a clause-density parameter; it stands as a benchmark
distribution, not a worst-case complexity class.  The name is needed to keep a
distributional phase-transition claim distinct from a worst-case complexity
claim, and is used here only to route the empirical literature
\cite{MitchellSelmanLevesque1992}.
\index{random k-SAT distribution}

\section{Decision diagrams and related problems}

\begin{hbdefinition}{def:p12-robdd}{Reduced ordered binary decision diagram}
Fix a finite variable set \(X\), a Boolean function
\(f:\{0,1\}^{X}\to\{0,1\}\), and a total order on \(X\).  An \emph{ordered
binary decision diagram} is a rooted acyclic decision graph in which variable
labels strictly increase along every root-to-terminal path in that order, so
no variable is tested twice on one path.  Each nonterminal node is labelled by
a variable and has exactly two ordered successors, its low child and high
child.  Given a total assignment, evaluation follows the low child when the
node's variable has value \(0\) and the high child when it has value \(1\).
The graph represents \(f\) exactly when, for every total assignment to \(X\),
the reached terminal's label equals the value of \(f\) on that assignment.
Terminal nodes are labelled by truth values, there is at most one
terminal carrying each truth-value label, and every node of the diagram is
reachable from its root.  A constant function is represented by the single
corresponding terminal root.
It is \emph{reduced} when no node has identical children
and no two distinct nodes have the same variable and the same ordered pair of
children.  The fixed order is necessary: without it the representation is not
canonical, and different orders can have exponentially different graph sizes.
Such a reduced ordered binary decision diagram is abbreviated \emph{ROBDD}.
The object is a data representation of a Boolean function, not a new logic;
its canonicity is used below, and its symbolic-model-checking application is
treated in Section~\ref{sec:p09-model-checking}.
\index{binary decision diagram}\index{reduced ordered binary decision diagram}
\index{ROBDD|see{reduced ordered binary decision diagram}}
\end{hbdefinition}

\begin{hbtheorem}{thm:p12-robdd-canonicity}{cited=Bryant1986}{ROBDD canonicity}
For every finite variable set \(X\), every total order on \(X\), and every
function \(f:\{0,1\}^{X}\to\{0,1\}\), there is, up to rooted labelled graph
isomorphism, exactly one reduced ordered binary decision diagram representing
\(f\).
\end{hbtheorem}

The quantified Boolean formula problem is already calibrated by
Theorem~\ref{thm:p04-tqbf-pspace}; the result is PSPACE-completeness for the
fixed encoding, not a statement about every bounded-alternation fragment.

\begin{hbdefinition}{def:p12-counting-and-maxsat}{Counting SAT and MaxSAT}
For a formula whose finite variable set is part of the input,
\(\#\mathsf{SAT}\) is the function returning the number of satisfying total
valuations.  It is a counting problem, not a decision language unless a
threshold question is added.  The class \(\#\mathsf P\) consists of functions
that return, for each input, the number of accepting branches of a
polynomial-time nondeterministic Turing machine.  This function-class standing
is needed to state the counting result and is not the decision class
\(\mathsf{NP}\).  For an unweighted finite clause multiset,
\(\mathsf{MaxSAT}\) asks for a valuation maximizing the number of satisfied
clauses; the threshold decision version asks whether at least \(k\) clauses
can be satisfied.  Weighted and partial variants must be declared separately.
For counting problems \(f\) and \(g\), a polynomial-time map \(h\) is a
\emph{parsimonious reduction} from \(f\) to \(g\) exactly when
\(f(x)=g(h(x))\) for every input \(x\).  It is a function-problem reduction
that preserves the exact witness count, not merely yes/no membership, and is
the reduction notion used by the counting theorem below.
These names are necessary to distinguish counting and optimization from
ordinary satisfiability, stand respectively for a function problem and an
optimization problem, and are used in the complexity statement below.
\index{counting SAT@\(\#\mathsf{SAT}\)}
\index{sharp-P class@\(\#\mathsf P\) class}
\index{parsimonious reduction}
\index{MaxSAT}
\end{hbdefinition}

\begin{hbtheorem}{thm:p12-counting-sat-status}{cited=Valiant1979Enumeration}{Counting-SAT complexity}
\(\#\mathsf{SAT}\) is complete for \(\#\mathsf P\) under polynomial-time
parsimonious reductions.
\end{hbtheorem}

\begin{hbtheorem}{thm:p12-maxsat-status}{cited=LiManya2009}{MaxSAT decision boundary}
The unweighted MaxSAT threshold problem is NP-complete:
membership uses a valuation certificate, and ordinary SAT is the special case
whose threshold is the total number of clauses.  The optimization variants
and their algorithms require their own encoding and approximation conventions
before a complexity or performance claim is made.
\end{hbtheorem}

\section{What the conflict-analysis apparatus does and does not supply}
\markright{59.4.\ CONFLICT-ANALYSIS APPARATUS}

\begin{hbdefinition}{def:p12-cdcl-apparatus}{Conflict-analysis vocabulary}
The name \emph{conflict-driven clause learning} (CDCL) is needed for the search
family that augments DPLL with learned entailed clauses and nonchronological
backtracking; it is not an alias for Method~\ref{def:p12-dpll}.  In the
apparatus below, an \emph{implication graph} has one node for each assigned
literal, annotated by decision level; an edge records use of an antecedent
clause in forcing an implied assignment; and a conflict node records a
falsified clause.  Decision level zero contains assignments made by initial
unit propagation before any branch decision; for \(k\geq 1\), level \(k\)
contains the \(k\)-th active branch decision and every assignment it propagates
before the next decision or a backtrack.  A \emph{unique implication point}
(UIP) is a current-level
node through which every path from the current decision to the conflict passes.
The \emph{first UIP} is the UIP closest to the conflict.  Resolving the
conflict explanation to a first-UIP clause yields an entailed
learned clause.  For a literal in that clause, its level means the decision
level of the currently assigned complementary literal that falsifies it.  A
learned clause is \emph{asserting} when exactly one of its literals is
falsified at the current decision level; that literal is the
\emph{asserting literal}, and the remaining clause literals are
\emph{nonasserting}.  The \emph{backjump level} is the greatest level among
the assigned complements of the nonasserting literals, or zero when there is
no such literal.  Backtracking to that level leaves the asserting literal
unassigned while the others remain false, so the learned clause is unit.
These clause roles are needed to determine this nonchronological return and
are not new object-language literal forms.  Watched literals are a
propagation data structure: two positions in a nonunit clause are monitored so
that the clause need not be revisited until a watched literal becomes false.
A restart discards the current partial assignment while retaining the input
and the selected learned clauses.  For a clause \(C\), reverse unit propagation
holds when unit propagation from the current database together with the
complements of every literal of \(C\) reaches the empty clause.  A clause has
the resolution-asymmetric-tautology property on a pivot \(L\in C\) when, for
every current clause containing \(\neg L\), the corresponding resolvent passes
that reverse-unit-propagation test.  A DRAT log is a sequence of clause
additions satisfying one of those checks and unrestricted deletions, ending
with addition of the empty clause; its independent checker validates the
refutation log, not the solver implementation.  These are metalinguistic
implementation and evidence roles, not new logical rules.  Their exact
distinctions are needed by
Table~\ref{tab:p12-dpll-cdcl} and Figure~\ref{fig:p12-cdcl-conflict}.
\index{conflict-driven clause learning}\index{implication graph}
\index{decision level}\index{unique implication point}\index{first UIP}
\index{asserting clause}\index{asserting literal}\index{backjump level}%
\index{nonasserting literal}\index{watched literal}\index{restart}
\index{reverse unit propagation}
\index{resolution-asymmetric tautology}
\index{DRAT}
\end{hbdefinition}

\begin{longtable}{@{}p{.16\textwidth}p{.34\textwidth}p{.40\textwidth}@{}}
\caption{DPLL and CDCL feature boundary.}
\label{tab:p12-dpll-cdcl}\\
\toprule
Feature & Displayed DPLL core & CDCL family \\ \midrule
\endfirsthead
\toprule
Feature & Displayed DPLL core & CDCL family \\ \midrule
\endhead
Propagation & Repeated unit propagation & Unit propagation, commonly
implemented with watched literals \\ \addlinespace
Branch failure & Chronological return to the most recent open branch &
Conflict analysis learns an entailed clause and may backjump
nonchronologically \\ \addlinespace
Learning & None & Typically first-UIP or another declared conflict-analysis
scheme; the scheme and deletion policy matter \\ \addlinespace
Restarts & None & May restart the assignment search while retaining selected
learned clauses; this changes search, not clause entailment \\ \addlinespace
Certificates & A returned valuation is directly checkable; the displayed
failure trace is an internal search result & DRAT or another declared proof log
can support independent UNSAT checking; a log does not verify the solver code \\
\addlinespace
Evidence here & Theorem~\ref{thm:p12-dpll-correctness}, proved in
Annex~\ref{annex:p12.sat.1} & Definitions and apparatus only \\
\bottomrule
\end{longtable}
The learning and propagation columns are sourced to
\cite{MarquesSilvaSakallah1999,MoskewiczEtAl2001} and
\cite{EenSorensson2003,MarquesSilvaLynceMalik2009}; the proof-log boundary is
sourced to \cite{WetzlerHeuleHunt2014}.

\begin{figure}[htbp]
\centering
\begin{tikzpicture}[
  >=Stealth,
  decision/.style={draw,rounded corners,fill=black!8,minimum width=1.15cm},
  implied/.style={draw,circle,minimum size=.72cm},
  conflict/.style={draw,rectangle,rounded corners,fill=black!15,minimum width=.9cm},
  every edge/.style={draw,->}
]
  \node[decision] (a) at (0,1.8) {\(a\;[1]\)};
  \node[implied] (b) at (2,1.8) {\(b\;[1]\)};
  \node[decision] (c) at (0,0) {\(c\;[2]\)};
  \node[implied] (d) at (2,0) {\(d\;[2]\)};
  \node[implied] (e) at (4,1.0) {\(e\;[2]\)};
  \node[implied] (f) at (4,-1.0) {\(f\;[2]\)};
  \node[implied] (g) at (6,1.0) {\(g\;[2]\)};
  \node[conflict] (x) at (7.6,0) {\(\bot\)};
  \path (a) edge node[above,sloped]{\scriptsize \(C_1\)} (b);
  \path (c) edge node[below,sloped]{\scriptsize \(C_2\)} (d);
  \path (b) edge node[above,sloped]{\scriptsize \(C_3\)} (e);
  \path (d) edge node[below,sloped]{\scriptsize \(C_3\)} (e);
  \path (d) edge node[below,sloped]{\scriptsize \(C_4\)} (f);
  \path (e) edge node[above,sloped]{\scriptsize \(C_5\)} (g);
  \path (f) edge node[below,sloped]{\scriptsize \(C_6\)} (x);
  \path (g) edge node[above,sloped]{\scriptsize \(C_6\)} (x);
\end{tikzpicture}
\caption{A conflict graph with learned clause and backjump.
The clauses are \(C_1=\neg a\lor b\), \(C_2=\neg c\lor d\),
\(C_3=\neg b\lor\neg d\lor e\), \(C_4=\neg d\lor f\),
\(C_5=\neg e\lor g\), and \(C_6=\neg f\lor\neg g\).
Square nodes are decisions, circles are implications, the filled box is the
conflict, bracketed numerals are
decision levels, and edge labels name antecedent clauses.  The node \(d\) is
the first UIP: every level-2 path from decision \(c\) to the conflict passes
through it.  Resolving the conflict explanation yields the entailed clause
\(\neg b\lor\neg d\); its nonasserting literal has level 1, so the search
backjumps to level 1 and the learned clause asserts \(\neg d\).}
\label{fig:p12-cdcl-conflict}
\end{figure}

\begin{hbdefinition}{def:p12-cdcl-transition}{CDCL transition state}
A CDCL state is \((D,\tau,\rho,k)\), where \(D\) is the input clause set plus
retained learned clauses, \(\tau\) is a consistent trail of signed assignments
partitioned into decision levels, \(\rho\) maps each propagated literal to its
antecedent clause, and \(k\) is the current level. The transitions are:
\begin{itemize}
  \item \emph{propagate}: append a literal forced by a unit clause and record
        that clause in \(\rho\);
  \item \emph{decide}: append an unassigned literal, incrementing \(k\);
  \item \emph{analyze}: when a clause is falsified, resolve it with antecedent
        clauses until an asserting entailed clause \(C\) is obtained;
  \item \emph{learn/backjump}: add \(C\), erase trail entries above its
        backjump level, and propagate its asserting literal;
  \item \emph{restart}: erase the nonzero trail while retaining the selected
        learned database;
  \item \emph{forget}: delete a learned clause not currently serving as a
        reason, subject to the declared certificate policy.
\end{itemize}
The state name is necessary because the database, trail, reasons, and level
have different invariants. It stands as an operational transition system,
differs from the DPLL method by learned clauses and nonchronological return,
and is used in the correctness and termination boundaries
\cite{MarquesSilvaSakallah1999,EenSorensson2003}.
\index{CDCL!transition system}
\end{hbdefinition}

\begin{hbtheorem}{thm:p12-cdcl-invariant}{cited=MarquesSilvaLynceMalik2009}{Learned-clause and backjump invariant}
At every transition reachable from the initial state, each learned clause in
\(D\) is entailed by the input clause set, each propagated trail literal is
justified by its recorded antecedent under the preceding trail, and a
first-UIP learned clause is asserting after the prescribed backjump. Hence
learning, backjumping, restarting, and permitted forgetting preserve the input
models; reaching the empty clause proves unsatisfiability, while a total
conflict-free trail satisfies the input.
\end{hbtheorem}

\noindent\emph{Proof.}
Propagation adds only a literal whose antecedent has every other literal false.
Decisions change the candidate assignment but not \(D\). Conflict analysis is
a sequence of resolution steps, so its learned clause is entailed by \(D\) and
therefore by the input, inductively. First-UIP analysis leaves exactly one
current-level literal; erasing higher levels leaves all other learned-clause
literals false and that literal unassigned, so unit propagation asserts it.
Restarts only erase assignments. Forgetting an entailed learned clause enlarges
the search state but does not change the input models. The terminal cases then
have the stated meanings.

If decisions are fair, forgetting is disabled, restarts do not prevent
progress, and every conflict learns a previously absent non-tautological
clause, only finitely many clauses over the fixed variable set can be learned,
so search terminates. Production CDCL permits deletion, restarts, and heuristic
choices and therefore needs a separate liveness policy for a termination
claim. VSIDS, phase saving, clause activity, and restart schedules alter search
order only; they carry no soundness premise. A DRAT/LRAT-style checker validates
the emitted refutation independently, and a model checker validates a SAT
assignment \cite{WetzlerHeuleHunt2014}.

The omega-automata/LTL procedure formerly owed here is now supplied in
Section~\ref{sec:p09-model-checking}, including acceptance conditions,
LTL-to-B\"uchi translation, and finite-state fixed-point algorithms. Symbolic
BDD model checking and its canonicity/order boundary are supplied there as
well. The ROBDD definition in this chapter remains the representation-level
interface and does not duplicate those Part-IX algorithms.
